\chapter{Testování a vyhodnocení}
\label{chap:evaluation}

V této kapitole ověříme, že knihovna PlanG zachovává chování původních
generátorů, a porovnáme její rychlost s nativními programy. Zajímá nás hlavně
to, jak velkou režii přidává C++ rozhraní a jak se tato režie projeví při
použití callbacků.

\section{Funkční testování}
\label{sec:evaluation-functional}

Nejdříve jsme ověřili, že wrapper pro malé instance vytváří stejný výstup jako
původní programy. Pro \texttt{geng} jsme porovnali výstup ve formátu
\texttt{graph6} pro grafy na šesti vrcholech. Pro \texttt{plantri} jsme stejným
způsobem porovnali triangulace na osmi vrcholech vypsané také ve formátu
\texttt{graph6}. V obou případech se soubory vytvořené přes PlanG shodovaly se
soubory vytvořenými nativním generátorem.

Dále jsme testovali běhy s callbacky \texttt{setPrune}, \texttt{setPreprune} a
\texttt{setOutproc}. Tyto testy ověřují, že se callbacky volají ve správné fázi
generování a že \texttt{GraphView} poskytuje aktuální graf bez nutnosti jeho
kopírování. Pro rozšíření \texttt{geng} o barevné třídy jsme testovali také
přesné velikosti barevných tříd, intervalové meze tříd a zakotvené vrcholy.

\section{Metodika měření}
\label{sec:evaluation-methodology}

Měření probíhalo na počítači s procesorem Intel Core i9-13900H v prostředí
Linux/WSL2. Knihovna byla sestavena překladačem \texttt{g++} 11.4.0 s
optimalizačními přepínači
\texttt{-O3 -flto -march=native -std=c++20}. Pro porovnání jsme použili
nativní binárky \texttt{geng} z balíku nauty 2.9.1 a \texttt{plantri} verze
5.5.

Nativní varianta \texttt{geng} a varianta integrovaná v PlanG byly sestaveny
se stejnými optimalizačními přepínači
\texttt{-O3 -flto -march=native -DNDEBUG}.

Aby měření nebylo ovlivněno rychlostí zápisu na výstup, vypnuli jsme vypisování
grafů. U nativních programů jsme použili přepínač \texttt{-u}; u \texttt{geng}
navíc přepínač \texttt{-q}, který potlačí pomocný výstup. V PlanG jsme použili
odpovídající metodu \texttt{setNoOutput}. Každý běh jsme spustili jako
samostatný proces a pomocí příkazu \texttt{time} jsme měřili jeho celkový čas
\texttt{real}. Jednotlivé konfigurace jsme spustili alespoň třikrát a v
tabulkách uvádíme medián naměřených časů.

\section{Čisté generování}
\label{sec:evaluation-baseline}

Nejprve porovnáme běh bez uživatelských callbacků. V tomto případě C++ část
knihovny pouze převede konfiguraci na argumenty původního programu, nastaví
prázdný runtime stav a spustí nativní generátor. Neočekáváme proto výraznou
ztrátu výkonu.

\begin{table}
\centering
\begin{tabular}{rrrrr}
\toprule
\(n\) & počet grafů & \texttt{geng} [s] & PlanG [s] & PlanG/\texttt{geng} \\
\midrule
8  & 12\,346           & 0{,}008   & 0{,}009   & 1{,}13 \\
9  & 274\,668          & 0{,}068   & 0{,}067   & 0{,}99 \\
10 & 12\,005\,168     & 2{,}101   & 2{,}169   & 1{,}03 \\
11 & 1\,018\,997\,864 & 237{,}542 & 243{,}459 & 1{,}02 \\
\bottomrule
\end{tabular}
\caption{Porovnání nativního \texttt{geng} a PlanG při generování bez výstupu.}
\label{tab:evaluation-geng-baseline}
\end{table}

Na základě výsledků uvedených v \cref{tab:evaluation-geng-baseline} můžeme
dojít k závěru, že při čistém generování dosahuje PlanG srovnatelného
výkonu jako nativní \texttt{geng}. V tomto scénáři tedy nepozorujeme
významnou režii způsobenou použitím rozhraní PlanG.

\begin{table}
\centering
\begin{tabular}{rrrrr}
\toprule
\(n\) & počet triangulací & \texttt{plantri} [s] & PlanG [s] & PlanG/\texttt{plantri} \\
\midrule
14 & 339\,722     & 0{,}125 & 0{,}123 & 0{,}98 \\
15 & 2\,406\,841  & 0{,}869 & 0{,}838 & 0{,}96 \\
16 & 17\,490\,241 & 6{,}169 & 6{,}080 & 0{,}99 \\
\bottomrule
\end{tabular}
\caption{Porovnání nativního \texttt{plantri} a PlanG při generování bez výstupu.}
\label{tab:evaluation-plantri-baseline}
\end{table}

Také u \texttt{plantri} zůstává běh přes PlanG srovnatelný s nativním programem,
jak ukazuje \cref{tab:evaluation-plantri-baseline}.

\section{Callbacky a srovnání s pluginem}
\label{sec:evaluation-callbacks}



\begin{table}
\centering
\begin{tabular}{rrrr}
\toprule
\(n\) & bez callbacku [s] & \texttt{setPrune} [s] & \texttt{setPreprune} [s] \\
\midrule
9  & 0{,}078 & 0{,}079 & 0{,}074 \\
10 & 2{,}408 & 2{,}417 & 2{,}463 \\
\bottomrule
\end{tabular}
\caption{Režie prázdných callbacků v backendu \texttt{geng}.}
\label{tab:evaluation-geng-callback-overhead}
\end{table}



\begin{table}
\centering
\begin{tabular}{rrrrr}
\toprule
\(n\) & počet triangulací & \texttt{plantri\_maxd} [s] & PlanG [s] & PlanG/native \\
\midrule
16 & 193\,445     & 0{,}295 & 0{,}306 & 1{,}04 \\
17 & 686\,364     & 1{,}120 & 1{,}157 & 1{,}03 \\
18 & 2\,444\,816 & 4{,}366 & 4{,}465 & 1{,}02 \\
\bottomrule
\end{tabular}
\caption{Porovnání nativního pluginu \texttt{maxdeg.c} a callbackové implementace v PlanG.}
\label{tab:evaluation-plantri-maxdeg}
\end{table}

TODO

\section{Cena bohatších struktur}
\label{sec:evaluation-colours}

TODO

\begin{table}
\centering
\begin{tabular}{rrrr}
\toprule
\(n\) & bez barev [s] & dvě vyvážené barvy [s] & dva zakotvené vrcholy [s] \\
\midrule
8 & 0{,}007 & 0{,}121 & 0{,}108 \\
9 & 0{,}070 & 4{,}958 & 3{,}498 \\
\bottomrule
\end{tabular}
\caption{Čas generování \texttt{geng} v PlanG při zapnutí barevných tříd.}
\label{tab:evaluation-coloured-geng}
\end{table}



\section{Shrnutí}
\label{sec:evaluation-summary}

TODO
